% =====================================================================
% V-C. Evaluación en distribución (test)  (E3)
% =====================================================================
\subsection{Evaluación en distribución (test)}
\label{sec:test}

El modelo elegido en valid (pérdida combinada, solo dorados) se evaluó una única vez en test con
sus tres semillas y umbral 0.5 (Tabla~\ref{tab:test}). Supera a Otsu en todas las métricas: el
IoU sube de 0.554 a 0.669 ($+0.114$). La mejora viene sobre todo de la precisión ($+0.133$); el
recall mejora menos ($+0.041$). Es decir, la red aprendió principalmente a \emph{no} confundir
con agua superficies oscuras que un umbral sí confunde.

\begin{table}[t]
\centering
\caption{Test (eventos vistos, recortes no vistos). U-Net: media $\pm$ desviación sobre 3
semillas, umbral 0.5.}
\label{tab:test}
\begin{tabular}{lcc}
\toprule
Métrica & Otsu & U-Net (BCE + Dice) \\
\midrule
IoU       & 0.554 & $\mathbf{0.669\pm0.003}$ \\
F1        & 0.713 & 0.801 \\
Precisión & 0.682 & 0.815 \\
Recall    & 0.748 & 0.789 \\
Omisión   & 0.252 & 0.211 \\
Comisión  & 0.318 & 0.185 \\
\bottomrule
\end{tabular}
\end{table}

\begin{figure}[t]
\centering
\includegraphics[width=0.85\columnwidth]{f7_matriz_confusion_test.png}
\caption{Matriz de confusión en test (suma de las 3 semillas).}
\label{fig:confusion}
\end{figure}

\textbf{Matriz de confusión (Fig.~\ref{fig:confusion}).} El error dominante es la
\emph{omisión}: el 21.1\,\% del agua real se predice como no agua. En cambio, solo el 2.5\,\% del
fondo se marca como agua. Esa tasa de comisión parece despreciable, pero como el fondo es mucho
más abundante que el agua, equivale a 1.35\,M de píxeles, suficientes para limitar la
precisión a 0.815. Leída por filas, la matriz dice que el modelo es confiable cuando afirma
``no agua'' y que, cuando falla, lo hace sobre todo por \emph{no ver} agua que sí está.

\textbf{Umbral.} La curva IoU--umbral es casi plana entre 0.2 y 0.8. El umbral óptimo en valid
($\approx$0.34) da 0.666 en test, sin diferencia apreciable con 0.5, así que mantenemos el
umbral fijado de antemano: las probabilidades de la red están bien separadas.

\textbf{Por evento.} El promedio esconde mucha variación: Nigeria alcanza 0.887, mientras
Pakistán cae a 0.194 con una omisión de 0.71. El agua perdida en Pakistán no es oscura en el
SAR: es el caso de agua no especular que analizamos en la Sección~\ref{sec:fallos}. La
Fig.~\ref{fig:errores_test} muestra mapas de error representativos.

\begin{figure}[t]
\centering
\includegraphics[width=\columnwidth]{f7_mapas_error_test.png}
\caption{Mapas de error en test. Azul: VP; rojo: FP (comisión); naranja: FN (omisión);
blanco: sin dato.}
\label{fig:errores_test}
\end{figure}
